\section{Experiments and Results}
\section{Experiments}
\label{sec:experiments}

We evaluate VersionGS from two perspectives:
whether versioning preserves the quality of the current map,
and whether historical maps can be stored efficiently and replayed
accurately.
We first compare current-scene rendering with existing Gaussian
mapping systems.
Historical fidelity, archive storage, and version loading are
evaluated separately.

\subsection{Experimental Setup}
\label{sec:experimental_setup}

\paragraph{Datasets.}
We evaluate on the synthetic Flat
dataset~\cite{schmid2022panoptic}
and the real-world Aria
dataset~\cite{pan2023ariadigitaltwinnew}.
Flat contains two RGB-D sequences of the same environment,
with substantial changes between captures.
For Aria, we use two rooms with two captures per room.
These sequences include object additions, removals, and rearrangements.
They allow us to assess adaptation to the latest observed scene
and preservation of earlier map versions.

\paragraph{Rendering metrics.}
We report PSNR, SSIM~\cite{wang2004image},
and LPIPS~\cite{zhang2018unreasonable}
for color rendering.
Depth accuracy is measured by mean absolute error,
reported as Depth L1 in centimeters.
Higher PSNR and SSIM indicate better quality,
while lower LPIPS and depth error indicate better quality.
We report input-view and novel-view results separately.
Rendering is evaluated at the dataset evaluation resolution
using the supplied camera poses.

\paragraph{Mapping baselines.}
We compare VersionGS with
SplaTAM~\cite{keetha2024splatam},
MonoGS~\cite{murai2024monogs},
DG-SLAM~\cite{xu2024dgslam},
and GaME~\cite{yugay2026game}.
SplaTAM and MonoGS provide established Gaussian mapping baselines,
while DG-SLAM addresses dynamic environments.
GaME is the mapping backbone of VersionGS and our main reference
for assessing how versioning affects current-scene quality.

\paragraph{Current-scene evaluation protocol.}
We follow the GaME protocol for current-scene rendering.
The captures of each scene are processed in temporal order
as a continuous RGB-D stream.
Every tenth frame of the final capture is reserved for
novel-view evaluation.
The remaining frames of that capture are available to the mapper
and are used for input-view evaluation.
Ground-truth camera poses are supplied to all methods
to isolate mapping performance from pose estimation.

VersionGS saves session commits and an additional commit after
final refinement.
The current-scene tables evaluate the final refined map restored
from its archive.
This evaluates the delivered renderable map rather than only
the model held in memory during training.
Session commits and their corresponding live maps are evaluated
separately because parameter reuse can introduce a small quality
difference during version construction.

\paragraph{Archive checking views.}
Version construction uses selected RGB-D observations to check
visibility and rendering quality.
These observations must be separate from the novel-view test set.
The final checking-view selection and data split will be reported
with the completed experiments.

\paragraph{Draft result status.}
The baseline values in
Tables~\ref{tab:rendering_flat_dataset}
and~\ref{tab:rendering_aria_dataset}
are reproduced from the supplied GaME comparison tables
and marked with \textsuperscript{R}.
VersionGS values marked with \textsuperscript{P}
are illustrative placeholders for manuscript preparation.
They are not measured results.
The corresponding numerical analysis is provisional
and must be replaced after evaluation.


\begin{figure*}[t]
    \centering
    \setlength{\vgsqualwidth}{0.151\textwidth}
    \setlength{\tabcolsep}{2pt}

    \begin{tabular}{
        @{}c
        *{6}{>{\centering\arraybackslash}m{\vgsqualwidth}}
        @{}
    }
        &
        \scriptsize SplaTAM &
        \scriptsize MonoGS &
        \scriptsize DG-SLAM &
        \scriptsize GaME &
        \scriptsize VersionGS &
        \scriptsize Ground Truth
        \\[3pt]

        \rotatebox{90}{\small A}
        &
        \vgsqualimage{
            figures/qualitative_comparisons/ADD_splatam__flat_dataset_run2_287_new.png
        }
        &
        \vgsqualimage{
            figures/qualitative_comparisons/ADD_monogs__flat_dataset_run2_286.png
        }
        &
        \vgsqualimage{
            figures/qualitative_comparisons/ADD_dgslam__flat_dataset_run2_286.png
        }
        &
        \vgsqualimage{
            figures/qualitative_comparisons/ADD_flat_dataset_run2_318.png
        }
        &
        \vgsqualimage{
            figures/qualitative_comparisons/versiongs_flat_add.png
        }
        &
        \vgsqualimage{
            figures/qualitative_comparisons/ADD_gt_flat_dataset_run2_318.png
        }
        \\[4pt]

        \rotatebox{90}{\small B}
        &
        \vgsqualimage{
            figures/qualitative_comparisons/aria_splatam_0_agent_3_frame_000200.png
        }
        &
        \vgsqualimage{
            figures/qualitative_comparisons/aria_monogs_0_agent_3_frame00020.png
        }
        &
        \vgsqualimage{
            figures/qualitative_comparisons/aria_dgslam_0_agent_3_frame200.png
        }
        &
        \vgsqualimage{
            figures/qualitative_comparisons/aria_render_0_agent_3_frame_000200.png
        }
        &
        \vgsqualimage{
            figures/qualitative_comparisons/versiongs_aria_rearrangement.png
        }
        &
        \vgsqualimage{
            figures/qualitative_comparisons/aria_gt_0_agent_3_frame_000200.jpg
        }
        \\[4pt]

        \rotatebox{90}{\small C}
        &
        \vgsqualimage{
            figures/qualitative_comparisons/remove_splatam_train_91_new.png
        }
        &
        \vgsqualimage{
            figures/qualitative_comparisons/remove_monogs_train_91.png
        }
        &
        \vgsqualimage{
            figures/qualitative_comparisons/remove_dgslam_train_91.png
        }
        &
        \vgsqualimage{
            figures/qualitative_comparisons/remove_ours_train_91.png
        }
        &
        \vgsqualimage{
            figures/qualitative_comparisons/versiongs_flat_rearrangement.png
        }
        &
        \vgsqualimage{
            figures/qualitative_comparisons/remove_gt_train_91.png
        }
        \\
    \end{tabular}

    \caption{
        \textbf{Current-scene rendering under scene changes.}
        The selected cases contain object additions and rearrangements.
        (A) A black office chair is added.
        (B) The toy house, chair, and picture are rearranged.
        (C) The cutlery is replaced, and the painting and chair are moved.
        GaME and VersionGS are shown separately.
        VersionGS panels are intended for renders from the restored
        final commit.
        Missing images are explicitly shown as placeholders.
        All panels in the completed comparison must use matched
        evaluation cameras.
    }
    \label{fig:qualitative_results}
\end{figure*}


\subsection{Evolving Scene Mapping Results}
\label{sec:evolving_mapping_results}

We first examine whether VersionGS preserves current-scene
rendering quality while adding version storage and replay.
Tables~\ref{tab:rendering_flat_dataset}
and~\ref{tab:rendering_aria_dataset}
compare input-view and novel-view rendering on Flat and Aria.
The comparison with GaME is particularly important because
both methods use the same mapping backbone.
It measures the effect of the additional parameter-management
rules on the final map.

All metrics are computed over complete images.
Changed objects often occupy only a small part of a scene,
so whole-image averages may hide errors near added, removed,
or rearranged objects.
We therefore complement the quantitative evaluation with
the selected regions in Fig.~\ref{fig:qualitative_results}.


\begin{table*}[t]
    \centering
    \small
    \setlength{\tabcolsep}{9pt}

    \resizebox{\textwidth}{!}{%
    \begin{tabular}{lcccc}
        \toprule
        \textbf{Method}
        & \textbf{PSNR [dB]} $\uparrow$
        & \textbf{SSIM} $\uparrow$
        & \textbf{LPIPS} $\downarrow$
        & \textbf{Depth L1 [cm]} $\downarrow$
        \\
        \midrule

        SplaTAM~\cite{keetha2024splatam}\vgsreported
        & 15.88 / 12.69
        & 0.484 / 0.299
        & 0.548 / 0.696
        & 21.16 / 59.91
        \\

        MonoGS~\cite{murai2024monogs}\vgsreported
        & 21.24 / 21.33
        & 0.771 / 0.769
        & 0.398 / 0.396
        & 30.95 / 29.91
        \\

        DG-SLAM~\cite{xu2024dgslam}\vgsreported
        & 13.72 / 13.70
        & 0.586 / 0.603
        & 0.737 / 0.739
        & 73.76 / 73.90
        \\

        GaME~\cite{yugay2026game}\vgsreported
        & 24.79 / 23.47
        & 0.899 / 0.880
        & 0.227 / 0.266
        & 16.63 / 15.59
        \\

        \midrule

        VersionGS (Ours)
        & \vgsplaceholder{25.00 / 23.70}
        & \vgsplaceholder{0.903 / 0.885}
        & \vgsplaceholder{0.220 / 0.258}
        & \vgsplaceholder{16.30 / 15.20}
        \\

        \bottomrule
    \end{tabular}%
    }

    \caption{
        \textbf{Current-scene rendering on Flat.}
        Cells report input / novel-view results.
        VersionGS is evaluated using the final refined map
        restored from the archive.
        \textsuperscript{R}: baseline values reproduced from
        the supplied GaME comparison table.
        \textsuperscript{P}: illustrative draft values,
        not measured results.
    }
    \label{tab:rendering_flat_dataset}
\end{table*}


\paragraph{Flat.}
\textit{Provisional analysis based on placeholder values.}
The draft values illustrate the intended comparison with GaME:
similar current-map quality while adding versioned storage.
They correspond to PSNR differences of
+0.21 dB on input views and +0.23 dB on novel views.
The placeholder SSIM values are also slightly higher,
while LPIPS and depth errors are lower.
These values are not experimental evidence.
If the trend is confirmed under a matched protocol and repeated runs,
it would support the conclusion that versioning preserves
current-scene reconstruction quality.
Ablation studies are needed to determine whether any improvement
comes from the freezing and unfreezing strategy or other
implementation differences.


\begin{table*}[t]
    \centering
    \small
    \setlength{\tabcolsep}{7pt}

    \resizebox{\textwidth}{!}{%
    \begin{tabular}{llcccc}
        \toprule
        \textbf{Method}
        & \textbf{Scene}
        & \textbf{PSNR [dB]} $\uparrow$
        & \textbf{SSIM} $\uparrow$
        & \textbf{LPIPS} $\downarrow$
        & \textbf{Depth L1 [cm]} $\downarrow$
        \\
        \midrule

        SplaTAM~\cite{keetha2024splatam}\vgsreported
        & room0
        & --\textsuperscript{$\dagger$} / 16.75
        & 0.780 / 0.447
        & 0.239 / 0.444
        & 4.12 / 17.79
        \\
        & room1
        & 22.94 / 17.90
        & 0.810 / 0.543
        & 0.219 / 0.440
        & 2.84 / 15.90
        \\
        & Avg.
        & 22.37 / 17.33
        & 0.795 / 0.495
        & 0.229 / 0.442
        & 3.48 / 16.80
        \\

        \midrule

        MonoGS~\cite{murai2024monogs}\vgsreported
        & room0
        & 25.28 / 25.19
        & 0.781 / 0.779
        & 0.277 / 0.273
        & 5.15 / 5.09
        \\
        & room1
        & 23.12 / 23.02
        & 0.844 / 0.842
        & 0.236 / 0.241
        & 4.99 / 5.01
        \\
        & Avg.
        & 24.20 / 24.11
        & 0.813 / 0.811
        & 0.257 / 0.257
        & 5.07 / 5.05
        \\

        \midrule

        DG-SLAM~\cite{xu2024dgslam}\vgsreported
        & room0
        & 15.78 / 15.63
        & 0.578 / 0.580
        & 0.761 / 0.763
        & 67.60 / 69.05
        \\
        & room1
        & 12.62 / 12.44
        & 0.651 / 0.643
        & 0.702 / 0.708
        & 55.47 / 57.03
        \\
        & Avg.
        & 14.20 / 14.04
        & 0.615 / 0.612
        & 0.732 / 0.736
        & 61.54 / 63.04
        \\

        \midrule

        GaME~\cite{yugay2026game}\vgsreported
        & room0
        & 28.87 / 28.94
        & 0.934 / 0.935
        & 0.166 / 0.166
        & 2.30 / 2.20
        \\
        & room1
        & 30.83 / 30.54
        & 0.961 / 0.958
        & 0.107 / 0.114
        & 1.80 / 1.90
        \\
        & Avg.
        & 29.85 / 29.74
        & 0.948 / 0.947
        & 0.137 / 0.140
        & 2.05 / 2.05
        \\

        \midrule

        VersionGS (Ours)
        & room0
        & \vgsplaceholder{29.10 / 29.20}
        & \vgsplaceholder{0.938 / 0.939}
        & \vgsplaceholder{0.160 / 0.160}
        & \vgsplaceholder{2.20 / 2.10}
        \\
        & room1
        & \vgsplaceholder{31.10 / 30.80}
        & \vgsplaceholder{0.965 / 0.962}
        & \vgsplaceholder{0.100 / 0.108}
        & \vgsplaceholder{1.70 / 1.80}
        \\
        & Avg.
        & \vgsplaceholder{30.10 / 30.00}
        & \vgsplaceholder{0.952 / 0.951}
        & \vgsplaceholder{0.130 / 0.134}
        & \vgsplaceholder{1.95 / 1.95}
        \\

        \bottomrule
    \end{tabular}%
    }

    \caption{
        \textbf{Current-scene rendering on Aria.}
        Cells report input / novel-view results.
        Avg.\ denotes the mean over the two rooms.
        \textsuperscript{R}: baseline values reproduced from
        the supplied GaME comparison table.
        \textsuperscript{P}: illustrative draft values,
        not measured results.
        $\dagger$: the supplied room0 input PSNR is inconsistent
        with the reported average and is omitted pending verification.
        The baseline average is retained as supplied.
    }
    \label{tab:rendering_aria_dataset}
\end{table*}


\paragraph{Aria.}
\textit{Provisional analysis based on placeholder values.}
The draft values show the same intended trend on real-world scenes.
Relative to the supplied GaME averages,
the placeholder PSNR differences are
+0.25 dB for input views and +0.26 dB for novel views.
The corresponding depth-error difference is -0.10 cm
for both view sets.
If confirmed experimentally, these results would indicate
that the versioning mechanisms remain compatible with accurate
mapping under real sensor noise and scene changes.
The final discussion will report measured differences and
variation across repeated runs.

\paragraph{Qualitative comparison.}
Figure~\ref{fig:qualitative_results}
organizes the comparison around object additions
and rearrangements.
The VersionGS panels must be generated by loading the final commit
in a separate replay process and rendering the same evaluation
cameras as the baselines.
This comparison checks whether the restored map correctly represents
the latest observed objects and their surroundings.
The analysis of visual differences will be completed
once the VersionGS renders are available.

\paragraph{Current quality versus historical preservation.}
Good rendering of the latest scene does not establish that
earlier versions have been preserved.
The current-scene tables therefore address only the compatibility
of versioning with the mapping backbone.
Historical rendering, exact parameter recovery, storage savings,
and loading costs require separate comparisons.
In particular, archival gains must be measured against
render-only snapshots with matched parameter precision
and compression settings.
